\section{Introduction}
\label{sec:introduction}

\subsection{Problem Statement}

The global proliferation of surveillance cameras has reached an unprecedented scale, with an estimated one billion CCTV cameras operating worldwide as of 2025~\cite{smith2025surveillance}. A single 1080p camera operating at 30 frames per second (fps) generates raw video data at approximately 1.5~Gbps, which even after hardware encoding typically produces a continuous stream of 2--20~Mbps depending on scene complexity and encoder configuration. For a mid-sized city deploying 10{,}000 cameras, this translates to 20--200~Gbps of aggregate video traffic and petabytes of storage consumed per month~\cite{wang2024urban}. These figures impose severe constraints on network bandwidth, storage infrastructure, and operational costs, often forcing operators to adopt aggressive compression at the expense of forensic image quality or to reduce camera frame rates during off-peak hours.

The problem is compounded by the fundamentally asymmetric nature of surveillance workloads. The vast majority of recorded footage---estimated at 95--99\% across typical deployments---contains no events of interest: empty parking lots at night, unoccupied lobbies, or static corridors~\cite{chen2023efficiency}. Yet conventional video codecs treat all frames uniformly, allocating bits according to a fixed rate-distortion (RD) model that makes no distinction between a motionless background and an active security incident. This uniform treatment wastes bandwidth on temporally redundant content while potentially under-allocating bits during critical moments when forensic detail is most needed.

\subsection{Motivation and Key Insight}

Our central insight is that surveillance video streams contain a strong, exploitable correlation between \emph{risk level} and \emph{information value}. When a scene is quiescent---no motion, no detected objects, no anomalous audio---the bits required to maintain forensic utility are minimal. Conversely, when risk indicators escalate (intrusion detection, abnormal motion patterns, audio anomalies), the value of each bit increases dramatically. A risk-aware compression system can exploit this asymmetry by dynamically allocating bits proportional to the instantaneous forensic importance of the scene, achieving substantial bandwidth savings during low-risk periods while preserving or even enhancing quality during critical events.

This paper presents an adaptive video compression framework that combines multi-modal risk scoring with formal rate-distortion optimization to achieve risk-proportional bandwidth allocation. Our system continuously evaluates visual, temporal, and audio signals to estimate the forensic risk of each frame, then uses a Lagrangian rate allocator to distribute a constrained bit budget across regions of interest (ROIs) and background regions. During low-risk periods, the system aggressively compresses static backgrounds---sometimes reducing them to as few as 120~kbps---while maintaining ROI quality sufficient for forensic identification. During high-risk periods, bit allocation shifts to maximize detail in forensically relevant regions, with pre- and post-event buffers ensuring complete evidence capture.

Our empirical evaluation on 1080p surveillance footage demonstrates that this approach achieves BD-rate improvements of up to 87.95\% over AV1 and 83.15\% over VP9, while operating at real-time throughput of 20.4~fps on edge hardware. The system preserves detection recall above 95\% across all risk states, ensuring that forensic evidence is never compromised by compression.

\subsection{Contributions}

This paper makes the following contributions:

\begin{enumerate}
    \item \textbf{Formal RD optimization with Lagrangian allocation.} We formulate risk-aware compression as a constrained optimization problem and derive a Lagrangian rate allocator that distributes bits across ROIs and background regions according to a risk-weighted distortion metric. The allocator operates online with $O(1)$ per-frame complexity, making it suitable for real-time edge deployment.

    \item \textbf{Multi-modal risk scoring.} We propose a risk scoring engine that fuses visual features (optical flow magnitude, CLAHE-enhanced edge density), temporal features (inter-frame change rate, scene transition detection), and optional audio features (ambient noise level, anomalous sound detection) into a single scalar risk score in $[0, 1]$. This score drives both the rate allocation and the adaptive GOP structure.

    \item \textbf{Forensic evidence preservation.} We introduce a buffer management scheme that maintains configurable pre-event and post-event quality windows around detected incidents. Frames within these windows are encoded at elevated quality regardless of instantaneous risk, ensuring that forensic evidence before and after a triggering event is preserved at investigation-grade fidelity.

    \item \textbf{Edge-cloud architecture for scalable deployment.} We describe an end-to-end system architecture that performs risk scoring and adaptive encoding at the network edge, with optional cloud-based storage and analytics. The architecture supports heterogeneous encoder backends (H.264, H.265, VP9, AV1) and degrades gracefully under bandwidth constraints through progressive quality reduction.
\end{enumerate}

\subsection{Paper Organization}

The remainder of this paper is organized as follows. Section~\ref{sec:related_work} reviews related work in video compression, ROI-based coding, adaptive bitrate streaming, and edge computing for video analytics. Section~\ref{sec:methodology} presents our system architecture, the risk scoring algorithm, the Lagrangian rate allocator, and the forensic evidence preservation mechanism. Section~\ref{sec:experiments} describes our experimental setup, evaluation metrics, baseline comparisons, and ablation study results. Section~\ref{sec:discussion} discusses findings, limitations, and future work directions. Section~\ref{sec:conclusion} concludes the paper.
